\documentclass[12pt]{article}

% =========================
% 0) Metadata (EDIT ME)
% =========================
\newcommand{\CourseCode}{MATH 521}
\newcommand{\CourseTitle}{Analysis I}
\newcommand{\CourseSection}{LEC 002}
\newcommand{\Semester}{Spring 2026}
\newcommand{\Instructor}{Thierry Laurens}
\newcommand{\MeetingInfo}{MWF 1:20--2:10 \;|\; VV B119} % optional
\newcommand{\HomeworkNumber}{9} % <-- change each week
\newcommand{\YourName}{Alejandro De La Torre}
\newcommand{\DueDate}{March 27, 2026} % <-- or set e.g. January 31, 2026

% =========================
% 1) Core math + proof tools
% =========================
\usepackage{amsmath, amssymb, amsthm}

% =========================
% 2) Lists and spacing
% =========================
\usepackage{enumitem}
\setlist[enumerate,1]{label=\textbf{(\alph*)}, leftmargin=2em, itemsep=0.25em}
\setlist[itemize]{leftmargin=1.5em, itemsep=0.25em}
\setlength{\parskip}{0.35em}
\setlength{\parindent}{0pt}

% =========================
% 3) Page geometry + header/footer
% =========================
\usepackage{geometry}
\geometry{margin=1in}

\usepackage{fancyhdr}
\pagestyle{fancy}
\fancyhf{}
\lhead{\CourseCode\ --- \CourseSection, \Semester \;(\MeetingInfo)}
\rhead{Homework \HomeworkNumber}
\cfoot{\thepage}
\renewcommand{\headrulewidth}{0.4pt}
% Fixes common fancyhdr warning about header height:
\setlength{\headheight}{15pt}
% Optional: reclaim a bit of vertical space after increasing headheight
\addtolength{\topmargin}{-2pt}

% =========================
% 4) Links (subtle)
% =========================
\usepackage[hidelinks]{hyperref}
\setcounter{tocdepth}{1}

% =========================
% 5) Quotes / figures
% =========================
\usepackage{csquotes}
\usepackage{graphicx}
\graphicspath{{images/}} % put figures in ./images/

% =========================
% 6) Simple scratch box (no tcolorbox)
% =========================
% A lightweight environment for brief planning / scratch work.
% This avoids any tcolorbox dependencies.
\newenvironment{scratch}{%
  \par\smallskip
  \noindent\textbf{Discussion \& Scratch Work}\begin{quote}\small
}{%
  \end{quote}\smallskip
}

% =========================
% 7) Lightweight theorem-like environments (optional)
% =========================
\newtheorem*{claim}{Claim}
\newtheorem*{lemma}{Lemma}
\newtheorem*{remark}{Remark}
\newtheorem{theorem}{Theorem}
\newtheorem{proposition}{Proposition}
\newtheorem{corollary}{Corollary}
\newtheorem{definition}{Definition}
\newtheorem{example}{Example}
\newtheorem{exercise}{Exercise}

% =========================
% 8) Problem header command (with TOC + PDF bookmarks)
% Usage: \problem[Optional short title]{<number>}
% =========================
\makeatletter
\newcommand{\problem}[2][]{%
  \def\prob@title{Problem #2\if\relax\detokenize{#1}\relax\else\ (\textit{#1})\fi}%
  \phantomsection%
  \pdfbookmark[1]{\prob@title}{prob#2}%
  \addcontentsline{toc}{section}{\prob@title}%
  \section*{\prob@title}%
  \vspace{-0.25em}\hrule\vspace{0.8em}%
}
\makeatother

% =========================
% 9) Title block
% =========================
\title{Homework \HomeworkNumber}
\author{\YourName \\
\CourseCode\ --- \CourseTitle \\
\CourseSection \\
Instructor: \Instructor}
\date{\DueDate}

\begin{document}
\maketitle

% ------------------ collaboration + sources ------------------
% \section*{Collaboration Statement}
% Write “I worked alone” or list collaborators / external resources.
% "I worked on this homework independently. No outside collaborators or resources were used."

\tableofcontents
\bigskip
\hrule
\bigskip

\newpage

% ============================================================
% Problems (duplicate blocks as needed)
% ============================================================

\problem[]{1}
Consider the metric space $(\mathbb{R}, d)$ with the usual metric
$d(x,y)=|x-y|$. For each of the following sets $A\subseteq \mathbb{R}$, find
its closure, isolated points, and limit points. Prove your answer.
\begin{enumerate}[label=\textbf{(\alph*)}]
  \item $A=[0,1)$
  \item $A=\mathbb{Z}$
  \item $A=\{x\in\mathbb{Q}: 0<x<1\}$
\end{enumerate}

\begin{scratch}
Before working through the sets, I want to keep the definitions straight.

Let $(X,d)$ be a metric space and let $A\subseteq X$.
A point $x\in X$ is an \emph{adherent point} of $A$ if every open ball around $x$ meets $A$:
\[
\forall r>0,\quad B_r(x)\cap A\neq\varnothing.
\]
In plain English: no matter how closely I zoom in around $x$, I still hit the set $A$. So adherent points are exactly the points that cannot be separated from $A$ by a small open ball.

The \emph{closure} of $A$ is the set of all adherent points:
\[
\overline{A}=\{x\in X:x\text{ is an adherent point of }A\}.
\]
So to identify $\overline{A}$, I should ask: which points have the property that every ball around them hits $A$?

An adherent point $x$ of $A$ is \emph{isolated} if there exists some radius $r>0$ such that
\[
B_r(x)\cap A=\{x\}.
\]
In plain English: $x$ belongs to $A$, but if I zoom in closely enough around $x$, it is the only point of $A$ that remains. So an isolated point is a point of $A$ that stands alone locally.

A point is a \emph{limit point} of $A$ if it is an adherent point but not isolated. Equivalently, $x\in A'$ if every ball around $x$ contains some point of $A$ different from $x$ itself. In plain English: no matter how closely I zoom in around $x$, I can always find other points of $A$ arbitrarily near it.

The key relation from lecture is
\[
\overline{A}=A\cup A'.
\]
So once I understand which points are limit points, the closure usually becomes much easier to identify.

\textbf{General strategy for each set.}
For each set $A$, I want to think in three passes.
\begin{enumerate}[label=\textbf{(\arabic*)}]
    \item First, guess which points are adherent by asking whether every ball around the point must hit $A$.
    \item Next, among those points, decide which are isolated and which are limit points by asking whether a sufficiently small ball can isolate the point from the rest of $A$.
    \item Finally, write down the closure either directly from the adherent-point definition or by using
    \[
    \overline{A}=A\cup A'.
    \]
\end{enumerate}

\textbf{(a) }$A=[0,1)$.
The first thing I notice is that every point already in $[0,1)$ should be adherent, because every ball around such a point contains the point itself and therefore meets $A$. The real question is whether there are extra adherent points outside $A$. The point $1$ is the most obvious candidate, because every ball around $1$ reaches slightly to the left and therefore meets $[0,1)$. On the other hand, points strictly less than $0$ or strictly greater than $1$ should fail to be adherent, because they have positive distance from the interval and I should be able to choose a ball missing $A$ completely.

For isolated points, I should ask whether any point of $[0,1)$ can ever be the only point of $A$ in a small ball. That seems unlikely, because every point of the interval has other points of the interval arbitrarily nearby. In particular, interior points clearly have nearby neighbors on both sides, and even the endpoint $0$ has points of $A$ arbitrarily close to it from the right. So I expect there to be no isolated points. That would mean every point of $[0,1]$ is actually a limit point, and then the closure should be $[0,1]$.

To prove this formally, I would show: every point of $[0,1]$ is adherent; every point outside $[0,1]$ is not adherent; no point of $[0,1)$ is isolated; and every point of $[0,1]$ is a limit point.

\textbf{(b) }$A=\mathbb{Z}$.
Again, every integer is automatically adherent because every ball around an integer contains that integer itself. The real question is whether non-integers are adherent. Here the picture is very different from intervals: if $x\notin\mathbb{Z}$, then there is some positive distance from $x$ to the nearest integer, so I should be able to choose a small ball around $x$ that contains no integers at all. That suggests the only adherent points are the integers themselves, so the closure should be $\mathbb{Z}$.

For isolated points, integers now look completely different from interval points: around a fixed integer $n$, I can choose a small enough radius, such as anything less than $\tfrac12$, so that the ball around $n$ contains no other integers. So each integer should be isolated. Once that is true, there should be no limit points at all, because a limit point must have other points of $A$ arbitrarily close by, and that never happens for integers.

So the expected picture is: adherent points are exactly $\mathbb{Z}$, every point of $\mathbb{Z}$ is isolated, the set of limit points is empty, and the closure is $\mathbb{Z}$.

To prove this, I would show: every integer is adherent; every non-integer admits a ball missing $\mathbb{Z}$; every integer admits a ball whose intersection with $\mathbb{Z}$ is just that single integer; and therefore there are no limit points.

\textbf{(c) }$A=\{x\in\mathbb{Q}:0<x<1\}$.
This set behaves very differently from part (b), because the rationals are dense in the reals. Every rational in $(0,1)$ is adherent, but now many points not in $A$ should also be adherent. In fact, any real number in $[0,1]$ should be adherent, because every ball around such a point contains rational numbers from $(0,1)$. That includes irrational points in $(0,1)$ as well as the endpoints $0$ and $1$. On the other hand, if a point lies strictly below $0$ or strictly above $1$, I should be able to choose a small ball disjoint from $(0,1)$, so those points should not be adherent. This suggests the closure should again be $[0,1]$.

For isolated points, density tells me the opposite of part (b): no point of $A$ should be isolated, because every ball around a rational point in $(0,1)$ contains infinitely many other rationals from $(0,1)$. In fact, every point of $[0,1]$ should be a limit point, because every ball around such a point contains points of $A$ different from the center. So I expect the isolated-point set to be empty and the limit-point set to be all of $[0,1]$.

To prove this formally, I would use density of $\mathbb{Q}$ in $\mathbb{R}$: show every point of $[0,1]$ is adherent by finding rationals from $(0,1)$ in every ball, show points outside $[0,1]$ are not adherent by separating them from the interval, and then show no point of $A$ is isolated because every ball around a point of $A$ contains other rational points of $(0,1)$.
\end{scratch}

\textbf{Proof.}\\
We work in $(\mathbb{R},d)$ with $d(x,y)=|x-y|$.

\medskip
\textbf{(a) $A=[0,1)$.}

\emph{Adherent points / closure.}
If $x\in[0,1)$, then for any $r>0$, $x\in B_r(x)\cap A$, so $x$ is adherent. 
If $x=1$, then for any $r>0$, $(1-r,1)\subset B_r(1)$ contains points of $[0,1)$, hence $1$ is adherent. 
If $x<0$, choose $r=\frac{|x|}{2}>0$; then $B_r(x)\subset(-\infty,0)$, so $B_r(x)\cap A=\varnothing$. 
If $x>1$, choose $r=\frac{x-1}{2}>0$; then $B_r(x)\subset(1,\infty)$, so $B_r(x)\cap A=\varnothing$. 
Thus $\overline{A}=[0,1]$.

\emph{Isolated points.}
Let $x\in[0,1)$. For any $r>0$, there exists $y\in(0,1)$ with $0<|y-x|<r$ (e.g.\ take $y=x+\min\{r/2,(1-x)/2\}$ if $x<1$, or $y=r/2$ if $x=0$). Hence $B_r(x)\cap A\neq\{x\}$, so no point is isolated.

\emph{Limit points.}
Since every $x\in[0,1]$ has points of $A$ arbitrarily close distinct from $x$, we have $A'=[0,1]$.

\medskip
\textbf{(b) $A=\mathbb{Z}$.}

\emph{Adherent points / closure.}
If $n\in\mathbb{Z}$, then $n\in B_r(n)\cap A$ for all $r>0$, so $n$ is adherent. 
If $x\notin\mathbb{Z}$, let $\delta=\min_{k\in\mathbb{Z}}|x-k|>0$ and take $r=\delta/2$; then $B_r(x)$ contains no integers, so $B_r(x)\cap A=\varnothing$. 
Thus $\overline{A}=\mathbb{Z}$.

\emph{Isolated points.}
For each $n\in\mathbb{Z}$, take $r=\tfrac12$; then $B_r(n)\cap\mathbb{Z}=\{n\}$, so every integer is isolated.

\emph{Limit points.}
No point has other integers arbitrarily close, hence $A'=\varnothing$.

\medskip
\textbf{(c) $A=\{x\in\mathbb{Q}:0<x<1\}$.}

\emph{Adherent points / closure.}
If $x\in[0,1]$, then by density of $\mathbb{Q}$ in $\mathbb{R}$, every ball $B_r(x)$ contains a rational $q$ with $0<q<1$, hence $B_r(x)\cap A\neq\varnothing$; so $[0,1]\subseteq\overline{A}$. 
If $x<0$, choose $r=\frac{|x|}{2}$ so $B_r(x)\subset(-\infty,0)$; if $x>1$, choose $r=\frac{x-1}{2}$ so $B_r(x)\subset(1,\infty)$; in either case $B_r(x)\cap A=\varnothing$. 
Thus $\overline{A}=[0,1]$.

\emph{Isolated points.}
Fix $x\in A$. For any $r>0$, by density of $\mathbb{Q}$ there exists $y\in\mathbb{Q}$ with $0<y<1$ and $0<|y-x|<r$; hence $B_r(x)\cap A\neq\{x\}$. So $A$ has no isolated points.

\emph{Limit points.}
Every $x\in[0,1]$ has points of $A$ arbitrarily close and distinct from $x$, hence $A'=[0,1]$.

\medskip
Therefore:
\[
\begin{aligned}
\text{(a)}\;& \overline{A}=[0,1],\quad A'=[0,1],\quad \text{no isolated points};\\
\text{(b)}\;& \overline{A}=\mathbb{Z},\quad A'=\varnothing,\quad \text{all points isolated};\\
\text{(c)}\;& \overline{A}=[0,1],\quad A'=[0,1],\quad \text{no isolated points}.
\end{aligned}
\]
\qed

\newpage

\problem[]{2}
Let $(X,d)$ be a metric space, let $A$ be a nonempty subset of $X$, and let $U$
be an open subset of $X$.
\begin{enumerate}[label=\textbf{(\alph*)}]
  \item Prove that $U\cap\overline{A}\subseteq \overline{U\cap A}$.
  \item Prove that $\overline{U}\cap\overline{A}=\overline{U\cap A}$. (Hint:
  Use part (a) and properties of closure from lecture.)
  \item Conclude that if $U\cap A=\varnothing$, then $U\cap\overline{A}=\varnothing$.
\end{enumerate}

\begin{scratch}
\end{scratch}

\textbf{Proof.}\\
\qed

\newpage

\problem[]{3}
Recall that $d_2$ is a metric on $\mathbb{R}^2$. Consider the subspace
\[
Y=\{(x,0): x\in\mathbb{R}\}
\]
with the induced metric $d_Y$.
\begin{enumerate}[label=\textbf{(\alph*)}]
  \item Provide an example of a set $A\subseteq Y$ where $A$ is open in
  $(Y,d_Y)$ but not open in $(\mathbb{R}^2,d_2)$.
  \item Prove that for any $B\subseteq Y$, if $B$ is closed in $(Y,d_Y)$ then
  $B$ is closed in $(\mathbb{R}^2,d_2)$.
\end{enumerate}

\begin{scratch}
\end{scratch}

\textbf{Proof.}\\
\qed

\newpage

\problem[]{4}
Let $(X,d)$ be a metric space, $\{x_n\}_{n\ge 1}$ a sequence in $X$, and
$x\in X$. Prove that the sequence $\{x_n\}_{n\ge 1}$ converges to $x$ if and
only if for every open set $U$ that contains $x$, we have $x_n\in U$ for all
but finitely many $n\in\mathbb{N}$.

\begin{scratch}
\end{scratch}

\textbf{Proof.}\\
\qed

\newpage

\problem[]{5}
Recall that $d_2$ is a metric on $\mathbb{R}^2$. Suppose that
$x^n=(x_1^n,x_2^n)$ is a sequence of points in $\mathbb{R}^2$.
\begin{enumerate}[label=\textbf{(\alph*)}]
  \item Prove that $|x_1^n-x_1|\le d_2(x^n,x)$ for any
  $x=(x_1,x_2)\in\mathbb{R}^2$.
  \item Prove that the sequence $\{x^n\}_{n\ge 1}$ converges in
  $(\mathbb{R}^2,d_2)$ if and only if the sequences $\{x_1^n\}_{n\ge 1}$ and
  $\{x_2^n\}_{n\ge 1}$ of real numbers both converge.
\end{enumerate}

\begin{scratch}
\end{scratch}

\textbf{Proof.}\\
\qed

\newpage

\problem[]{6}
Let $(X,d)$ be a metric space and $\{x_n\}_{n\ge 1}$ a sequence in $X$. Prove
that if $\{x_n\}_{n\ge 1}$ converges to $x\in X$, then for any $y\in X$, the
sequence $\{d(x_n,y)\}_{n\ge 1}$ of real numbers converges to $d(x,y)$.

\begin{scratch}
\end{scratch}

\textbf{Proof.}\\
\qed

% ============================================================
% Optional appendix: keep a personal toolbox of definitions/theorems
% ============================================================
\newpage
\appendix
\section*{Appendix: Toolbox}
\addcontentsline{toc}{section}{Appendix: Toolbox}

\subsection*{Lecture 22}

\begin{definition}
Let $(X,d)$ be a metric space and $A\subseteq X$. We say:
\begin{enumerate}[label=\textbf{(\arabic*)}]
  \item $a\in X$ is an interior point of $A$ if $\exists\, r>0$ s.t.
  $B_r(a)\subseteq A$.
  \item
  \[
  A^\circ=\{a\in X:\ a \text{ is an interior point of } A\}
  \]
  is the interior of $A$.
  \item $A$ is open if $A^\circ=A$.
\end{enumerate}
\end{definition}

\begin{example}
$X=\mathbb{R}^2$ with the metric $d_2$, and
$A=[0,1]\times[0,1]$.
\begin{itemize}
  \item $x=\left(\frac{3}{4},\frac{1}{2}\right)$ is an interior point of $A$:
  $B_{1/4}(x)\subseteq A$.
  \item $y=\left(\frac{1}{2},1\right)$ is not an interior point of $A$:
  $\forall r>0,\ B_r(y)\nsubseteq A$.
  \item $A^\circ=(0,1)\times(0,1)$.
  \item $A$ is not open in $(\mathbb{R}^2,d_2)$.
\end{itemize}
\end{example}

\begin{example}
Let $(X,d)$ be a metric space.
\begin{itemize}
  \item $\varnothing$ and $X$ are open sets.
  \item $B_r(a)$ is an open set $\ \forall a\in X$ and $\forall r>0$:

  Claim: $B_r(a)\subseteq (B_r(a))^\circ$.

  Fix $x\in B_r(a)$, i.e.\ $d(x,a)<r$.

  Want: $x\in (B_r(a))^\circ$, i.e.\ $\exists\, s>0$ s.t.
  $B_s(x)\subseteq B_r(a)$. Set
  \[
  s=r-d(x,a)>0.
  \]
  Then
  \[
  y\in B_s(x)\Rightarrow d(y,x)<s
  \Rightarrow d(y,a)\le d(y,x)+d(x,a)<s+d(x,a)=r.
  \]
  So $B_s(x)\subseteq B_r(a)$.
\end{itemize}
\end{example}

\begin{remark}
Note that $A^\circ\subseteq A$ for any set $A$:
\[
x\in A^\circ \Rightarrow \exists\, r>0 \text{ s.t. } B_r(x)\subseteq A
\Rightarrow x\in B_r(x)\subseteq A
\]
since $d(x,x)=0<r$.

So, to prove that a set is open, it suffices to show $A^\circ\supseteq A$.
\end{remark}

\begin{proposition}
Let $(X,d)$ be a metric space and $A,B\subseteq X$.
\begin{enumerate}[label=\textbf{(\arabic*)}]
  \item If $A\subseteq B$, then $A^\circ\subseteq B^\circ$.
  \item $A^\circ\cup B^\circ\subseteq (A\cup B)^\circ$.
  \item $A^\circ\cap B^\circ=(A\cap B)^\circ$.
  \item $(A^\circ)^\circ=A^\circ$. In particular, $A^\circ$ is an open set.
  \item $A^\circ$ is the largest open set contained in $A$.
  \item A finite intersection of open sets is open.
  \item An arbitrary union of open sets is open.
\end{enumerate}
\end{proposition}

\begin{remark}
An arbitrary intersection of open sets need not be open. Eg.\ in
$(\mathbb{R},|\cdot|)$,
\[
\bigcap_{n\in\mathbb{N}}\left(-\frac{1}{n},\frac{1}{n}\right)=\{0\}.
\]
\end{remark}

\subsection*{Lecture 23}

\begin{definition}
Let $(X,d)$ be a metric space. A set $A\subseteq X$ is closed if $A^c$ is open.
\end{definition}

\begin{example}
Let $(X,d)$ be a metric space.
\begin{enumerate}[label=\textbf{(\arabic*)}]
  \item $\varnothing$ and $X$ are closed: since $\varnothing^c=X$ and
  $X^c=\varnothing$ are open.
  \item
  \[
  (B_r(a))^c=\{x\in X:\ d(x,a)\ge r\}
  \]
  is closed $\ \forall r>0$ and $\forall a\in X$: since
  \[
  [(B_r(a))^c]^c=B_r(a)
  \]
  is open.
  \item
  \[
  A=\{x\in X:\ d(x,a)\le r\}
  \]
  is closed $\ \forall r>0$ and $\forall a\in X$.

  Want: $A^c=\{x\in X:\ d(x,a)>r\}$ is open. Fix $x\in A^c$. Then $d(x,a)>r$.
  Set $s=d(x,a)-r>0$. Then $B_s(x)\subseteq A^c$ since
  \[
  y\in B_s(x)\Rightarrow d(y,x)<s
  \Rightarrow d(a,x)\le d(a,y)+d(y,x)<d(a,y)+s
  \Rightarrow d(a,y)>d(a,x)-s=r.
  \]
\end{enumerate}
\end{example}

\begin{remark}
``closed'' $\neq$ ``not open''. Eg., in $(\mathbb{R},|\cdot|)$,
\begin{itemize}
  \item $(0,1)$ is open
  \item $[0,1]$ is closed
  \item $\varnothing$ and $\mathbb{R}$ are open and closed
  \item $(0,1]$ is not open and not closed
\end{itemize}
\end{remark}

\begin{definition}
Let $(X,d)$ be a metric space and $A\subseteq X$. We say:
\begin{enumerate}[label=\textbf{(\arabic*)}]
  \item $x\in X$ is an adherent point of $A$ if:
  \[
  \forall r>0,\ B_r(x)\cap A\neq\varnothing.
  \]
  \item
  \[
  \overline{A}=\{x\in X:\ x \text{ is an adherent point of } A\}
  \]
  is the closure of $A$.
  \item An adherent point $x$ of $A$ is called:
  \begin{itemize}
    \item isolated if $\exists\, r>0$ s.t. $B_r(x)\cap A=\{x\}$
    \item a limit point of $A$ otherwise
  \end{itemize}
  \item
  \[
  A'=\{x\in X:\ x \text{ is a limit point of } A\}.
  \]
\end{enumerate}
\end{definition}

\begin{remark}
$A\subseteq \overline{A}$ for any set $A$:
\[
x\in A \Rightarrow x\in B_r(x)\cap A \qquad \forall r>0.
\]
\[
x\in A' \Leftrightarrow \text{ the intersection of } B_r(x) \text{ and }
A\setminus\{x\} \text{ is nonempty } \forall r>0
\]
\[
\overline{A}=A\cup A'.
\]
\end{remark}

\begin{example}
$(\mathbb{R},|\cdot|),\qquad A=\left\{\frac{1}{n}: n\in\mathbb{N}\right\}$.
\begin{enumerate}[label=\textbf{(\arabic*)}]
  \item $\frac{1}{n}$ is an isolated point of $A$ $\ \forall n\in\mathbb{N}$:

  Fix $n$. Set
  \[
  r=\frac{1}{n}-\frac{1}{n+1}>0.
  \]
  Then
  \[
  B_r\left(\frac{1}{n}\right)\cap A=\left\{\frac{1}{n}\right\}.
  \]
  \item $0$ is a limit point for $A$:

  Fix $r>0$. Pick $n\in\mathbb{N}$ s.t. $\frac{1}{n}<r$.
  Then
  \[
  \frac{1}{n}\in (-r,r)\cap A,
  \]
  so $(-r,r)\cap A\neq\varnothing$.
  \item
  \[
  \overline{A}=A\cup\{0\}.
  \]
\end{enumerate}
\end{example}

\begin{lemma}
Let $(X,d)$ be a metric space and $A\subseteq X$.
\begin{enumerate}[label=\textbf{(\arabic*)}]
  \item $(\overline{A})^c=(A^c)^\circ$
  \item $(A^\circ)^c=\overline{A^c}$
  \item $A=\overline{A}\iff A$ is a closed set.
\end{enumerate}
\end{lemma}

\begin{exercise}
Let $(X,d)$ be a metric space and $A,B\subseteq X$.
\begin{enumerate}[label=\textbf{(\arabic*)}]
  \item If $A\subseteq B$ then $\overline{A}\subseteq \overline{B}$.
  \item $\overline{A\cap B}\subseteq \overline{A}\cap\overline{B}$.
  \item $\overline{A\cup B}=\overline{A}\cup\overline{B}$.
  \item $\overline{\overline{A}}=\overline{A}$. In particular,
  $\overline{A}$ is a closed set.
  \item $\overline{A}$ is the smallest closed set containing $A$.
  \item A finite union of closed sets is a closed set.
  \item An arbitrary intersection of closed sets is a closed set.
\end{enumerate}

Use the proposition from last lecture and take the complement of each statement.

Eg., for \textbf{(1)},
\[
A\subseteq B \Rightarrow A^c\supseteq B^c \Rightarrow (A^c)^\circ\supseteq (B^c)^\circ
\Rightarrow \overline{A}^c\supseteq \overline{B}^c \Rightarrow \overline{A}\subseteq \overline{B}.
\]
\end{exercise}

\begin{definition}
Let $(X,d_X)$ be a metric space and let $\varnothing\neq Y\subseteq X$. Then
$d_Y:Y\times Y\to\mathbb{R}$,
\[
d_Y(x,y)=d_X(x,y)
\]
is a metric on $Y$, called the induced metric on $Y$. The metric space
$(Y,d_Y)$ is called a subspace of $(X,d_X)$.
\end{definition}

\begin{example}
$\mathbb{R}$ with $d_{\mathbb{R}}(x,y)=|x-y|$ is a metric space.

$[0,1]$ with $d_{[0,1]}(x,y)=|x-y|$ is a subspace.
\end{example}

\subsection*{Lecture 24}

\begin{proposition}
Let $(X,d_X)$ be a metric space and $(Y,d_Y)$ a subspace (with the induced
metric).
\begin{enumerate}[label=\textbf{(\arabic*)}]
  \item A set $V\subseteq Y$ is open in $(Y,d_Y)\ \Leftrightarrow\ \exists\, U\subseteq X$
  open in $(X,d_X)$ s.t.\ $V=U\cap Y$.
  \item A set $G\subseteq Y$ is closed in $(Y,d_Y)\ \Leftrightarrow\ \exists\, F\subseteq X$
  closed in $(X,d_X)$ s.t.\ $G=F\cap Y$.
\end{enumerate}
\end{proposition}

\begin{example}
\begin{enumerate}[label=\textbf{(\arabic*)}]
  \item $[0,1)$ is not open in $(\mathbb{R},|\cdot|)$.

  $[0,1)$ is open in $([0,2],|\cdot|)$, since
  \[
  [0,1)=(-1,1)\cap[0,2].
  \]
  \item $[0,1)$ is not closed in $(\mathbb{R},|\cdot|)$.

  $[0,1)$ is closed in $((-1,1),|\cdot|)$, since
  \[
  [0,1)=(-1,1)\cap[0,2].
  \]
\end{enumerate}
\end{example}

\begin{corollary}
Let $(X,d_X)$ be a metric space and $(Y,d_Y)$ a subspace. The following are
equivalent:
\begin{enumerate}[label=\textbf{(\arabic*)}]
  \item $Y$ is open (closed) in $X$.
  \item Any $A\subseteq Y$ that is open (closed) in $Y$ is also open (closed)
  in $X$.
\end{enumerate}
\end{corollary}

\begin{proposition}
Let $(X,d_X)$ be a metric space and $(Y,d_Y)$ a subspace. For any $A\subseteq Y$,
\[
\overline{A}^{\,Y}=\overline{A}^{\,X}\cap Y.
\]
\end{proposition}

\begin{definition}
Let $(X,d)$ be a metric space and let $\{x_n\}_{n\ge 1}$ be a sequence in $X$.
We say $\{x_n\}_{n\ge 1}$ converges to $x\in X$ if:
\[
\forall \varepsilon>0\ \exists N\in\mathbb{N}\ \text{ s.t. }\ d(x_n,x)<\varepsilon
\qquad \forall n\ge N.
\]
When this is true, we write ``$\lim_{n\to\infty} x_n=x$'' or
``$x_n\to x$ in $(X,d)$ as $n\to\infty$''.
\end{definition}

\begin{example}
\begin{enumerate}[label=\textbf{(\arabic*)}]
  \item For $(\mathbb{R},|\cdot|)$, we have
  \[
  d(x_n,x)=|x_n-x|.
  \]
  This agrees with the definition from lecture 7.
  \item For $(\mathbb{R}^2,d_p)$, we have
  \[
  d(x_n,x)=\left(|x_{n,1}-x_1|^p+|x_{n,2}-x_2|^p\right)^{1/p}.
  \]
\end{enumerate}
\end{example}

\begin{remark}
$x_n\to x$ in $(X,d)$ as $n\to\infty \ \Leftrightarrow\ $ the sequence of real
numbers $d(x_n,x)$ converges to $0$ as $n\to\infty$.
\end{remark}

\begin{lemma}
The limit of a sequence is unique, if it exists.
\end{lemma}

\begin{theorem}
A sequence $\{x_n\}_{n\ge 1}$ converges to $x \ \Leftrightarrow\ $ every
subsequence $\{x_{n_k}\}_{k\ge 1}$ converges to $x$.
\end{theorem}

\begin{definition}
Let $(X,d)$ be a metric space. We say that a sequence $\{x_n\}_{n\ge 1}$ is
Cauchy if:
\[
\forall \varepsilon>0\ \exists N\in\mathbb{N}\ \text{ s.t. }\ d(x_n,x_m)<\varepsilon
\qquad \forall n,m\ge N.
\]
\end{definition}

\begin{proposition}
If $\{x_n\}_{n\ge 1}$ converges, then it is Cauchy.
\end{proposition}

\begin{remark}
Not every Cauchy sequence converges! Eg.,
\begin{enumerate}[label=\textbf{(\arabic*)}]
  \item $((0,1),|\cdot|)$ is a metric space.

  $\left\{\frac{1}{n}\right\}_{n\ge 1}$ is Cauchy, but does not converge in
  $(0,1)$.
  \item $(\mathbb{Q},|\cdot|)$ is a metric space.
  \[
  x_1=3,\qquad x_{n+1}=\frac{x_n}{2}+\frac{3}{2x_n}\qquad \forall n\ge 1.
  \]
  Then $\{x_n\}_{n\ge 1}$ is Cauchy, but doesn't converge in $\mathbb{Q}$.

  By lecture 12, this sequence converges to $\sqrt{3}$.
\end{enumerate}
\end{remark}

\subsection*{Lecture 25}

\begin{definition}
A metric space $(X,d)$ is complete if every Cauchy sequence in $X$ converges in
$X$.
\end{definition}

\begin{example}
\begin{enumerate}[label=\textbf{(\arabic*)}]
  \item $(\mathbb{R},|\cdot|)$ is complete (by lecture 12)
  \item $(\mathbb{Q},|\cdot|)$ is not complete (by the previous example)
\end{enumerate}
\end{example}

\begin{proposition}
Let $1\le p\le\infty$. $\ \forall n\in\mathbb{N},\ (\mathbb{R}^n,d_p)$ is
complete.
\end{proposition}

\begin{proposition}
Let $(X,d)$ be a metric space, $A\subseteq X$, and $x\in X$. Then:
\[
x\in\overline{A}\ \Leftrightarrow\ \exists\ \text{a sequence }\{a_n\}_{n\ge 1}
\text{ in } A \text{ that converges to } x.
\]
\end{proposition}

\begin{remark}
Recall:
\[
\overline{A}=A\cup A'
\]
\[
x\in A' \Leftrightarrow \exists\ \text{a sequence in } A\setminus\{x\}
\text{ that converges to } x.
\]
\end{remark}

\begin{definition}
Let $(X,d)$ be a metric space and $A\subseteq X$. We say:
\begin{enumerate}[label=\textbf{(\arabic*)}]
  \item a collection of sets $\{U_i\}_{i\in I}$ is an open cover of $A$, if
  $I\neq\varnothing$ is a set, $U_i$ is open $\ \forall i\in I$, and
  \[
  A\subseteq \bigcup_{i\in I} U_i.
  \]
  \item the set $A$ is compact if: $\ \forall\ $ open cover
  $\{U_i\}_{i\in I}$ of $A$, $\exists\ $ a finite set $J\subseteq I$ s.t.
  $\{U_j\}_{j\in J}$ is an open cover of $A$. We call $\{U_j\}_{j\in J}$ a
  finite subcover.
\end{enumerate}
\end{definition}

\begin{example}
Any finite set $\{x_1,\ldots,x_n\}\subseteq X$ is compact.

Let $\{U_i\}_{i\in I}$ be an open cover of $\{x_1,\ldots,x_n\}$.

Then
\[
x_1\in \bigcup_{i\in I} U_i \Rightarrow \exists\, i_1\in I \text{ s.t. } x_1\in U_{i_1}
\]
\[
x_2\in \bigcup_{i\in I} U_i \Rightarrow \exists\, i_2\in I \text{ s.t. } x_2\in U_{i_2}
\]
\[
\vdots
\]

So $\{U_{i_1},\ldots,U_{i_n}\}$ is a finite subcover.
\end{example}

\end{document}
